\documentclass[11pt,a4paper]{article}
\usepackage[utf8]{inputenc}
\usepackage[T1]{fontenc}
\usepackage[french]{babel}
\usepackage{geometry}
\usepackage{hyperref}
\usepackage{graphicx}
\usepackage{listings}
\usepackage{xcolor}
\usepackage{booktabs}

\geometry{margin=2.5cm}

\hypersetup{
    colorlinks=true,
    linkcolor=blue,
    urlcolor=blue,
    citecolor=blue
}

\lstset{
    basicstyle=\ttfamily\small,
    breaklines=true,
    frame=single,
    backgroundcolor=\color{gray!10},
    keywordstyle=\color{blue},
    stringstyle=\color{red},
    commentstyle=\color{green!60!black}
}

\title{\textbf{Rapport Technique}\\Service Reviews - FlexShop}
\author{FlexShop Team}
\date{\today}

\begin{document}

\maketitle

\section{Présentation du projet}

Ce document présente le service \texttt{reviews-service}, un microservice Symfony dédié à la gestion des avis clients pour la marketplace FlexShop.

\subsection{Lien vers le dépôt Git}

Le code source complet est disponible sur GitHub :

\begin{center}
\url{https://github.com/GnsYnns/reviews-service}
\end{center}

\section{Architecture et endpoints}

\subsection{Endpoints API}

\begin{table}[h]
\centering
\begin{tabular}{lll}
\toprule
\textbf{Méthode} & \textbf{Endpoint} & \textbf{Description} \\
\midrule
GET & /api/reviews & Liste paginée des avis \\
GET & /api/reviews/\{id\} & Détail d'un avis \\
POST & /api/reviews & Créer un avis \\
PUT & /api/reviews/\{id\} & Modifier un avis \\
DELETE & /api/reviews/\{id\} & Supprimer un avis \\
\midrule
GET & /api/stats/reviews/average/\{productId\} & Moyenne des notes \\
GET & /api/stats/reviews/search?q=... & Recherche full-text \\
\bottomrule
\end{tabular}
\caption{Endpoints du service Reviews}
\end{table}

\subsection{Swagger UI}

La documentation interactive est accessible via Swagger UI :

\begin{center}
\url{http://localhost:8000/api}
\end{center}


\section{Difficultés rencontrées et solutions}

\subsection{Difficulté 1 : Conflit de port avec le service Catalogue}

\textbf{Problème :} Le service Symfony (port 8000) interceptait les appels vers le Catalogue Django configuré sur le même port.

\textbf{Solution :} Modification de l'URL du Catalogue dans \texttt{ProductExistsValidator.php} :

\begin{lstlisting}[language=PHP]
private string $catalogUrl = 'http://localhost:8001/api/v1'
\end{lstlisting}

Le Catalogue Django doit maintenant tourner sur le port 8001.

\subsection{Difficulté 2 : Conflit de routes API Platform}

\textbf{Problème :} Les endpoints custom \texttt{/api/reviews/search} et \texttt{/api/reviews/average} entraient en conflit avec le routage d'API Platform, retournant des erreurs 404.

\textbf{Solution :} Changement du préfixe de route pour éviter le conflit :

\begin{lstlisting}[language=PHP]
#[Route('/api/stats')]
class ReviewStatsController extends AbstractController
\end{lstlisting}

Les endpoints sont maintenant accessibles via \texttt{/api/stats/reviews/...}.

\subsection{Difficulté 3 : Content-Type pour API Platform}

\textbf{Problème :} API Platform rejette les requêtes avec \texttt{Content-Type: application/json}, exigeant \texttt{application/ld+json}.

\textbf{Solution :} Mise à jour des tests PHPUnit et documentation pour utiliser le bon header :

\begin{lstlisting}
Content-Type: application/ld+json
\end{lstlisting}

\subsection{Difficulté 4 : Configuration base de données test}

\textbf{Problème :} PHPUnit tentait de se connecter à une base \texttt{reviews\_test} qui n'existait pas (suffixe automatique Doctrine).

\textbf{Solution :} Suppression de la configuration \texttt{dbname\_suffix} dans \texttt{config/packages/doctrine.yaml} :

\begin{lstlisting}[language=YAML]
when@test:
    doctrine:
        dbal:
            # dbname_suffix supprime
\end{lstlisting}

\subsection{Difficulté 5 : Cache HTTP pour le Catalogue}

\textbf{Problème :} Chaque validation de produit déclenchait un appel HTTP vers le Catalogue, risquant de surcharger le service.

\textbf{Solution :} Implémentation d'un cache de 60 secondes :

\begin{lstlisting}[language=PHP]
$exists = $this->cache->get($cacheKey, function () use ($value) {
    $response = $this->httpClient->request(
        'GET',
        $this->catalogUrl . '/products/' . $value . '/',
        ['timeout' => 3]
    );
    return $response->getStatusCode() !== 404;
});
\end{lstlisting}

Pool de cache dédié dans \texttt{cache.yaml} :

\begin{lstlisting}[language=YAML]
catalog.cache:
    adapter: cache.app
    default_lifetime: 60
\end{lstlisting}

\section{Collection Postman}

Une collection Postman exportée est disponible dans le dossier \texttt{docs/} :

\begin{verbatim}
Fichier : docs/FlexShop-Reviews-API.postman_collection.json
\end{verbatim}

Cette collection contient les requêtes suivantes :
\begin{itemize}
    \item GET /api/reviews (liste paginée)
    \item GET /api/reviews?productId=1 (filtrage)
    \item POST /api/reviews (création avec validation)
    \item GET /api/stats/reviews/average/1 (moyenne)
    \item GET /api/stats/reviews/search?q=excellent (recherche)
\end{itemize}

\section{Validation et tests}

\subsection{Tests PHPUnit}

Les tests sont exécutés avec :

\begin{lstlisting}
php bin/phpunit
\end{lstlisting}

Résultats des 3 tests implémentés :
\begin{itemize}
    \item \texttt{testGetReviewsReturnsOk} - Vérifie le endpoint liste
    \item \texttt{testGetNonExistentReviewReturns404} - Vérifie la gestion des 404
    \item \texttt{testCreateReviewWithInvalidRatingFails} - Vérifie la validation (rating 1-5)
\end{itemize}

\subsection{Points d'attention implémentés}

\begin{itemize}
    \item \textbf{Timeout HTTP} : 3 secondes sur les appels au Catalogue
    \item \textbf{Gestion Catalogue DOWN} : Try/catch avec message "Catalogue indisponible"
    \item \textbf{Validation custom} : \texttt{ProductExists} vérifie l'existence avant création
    \item \textbf{Cache Catalogue} : 60 secondes pour limiter les appels réseau
\end{itemize}

\section{Conclusion}

Le service Reviews est fonctionnel avec :
\begin{itemize}
    \item API CRUD complète via API Platform
    \item Endpoints custom pour statistiques et recherche
    \item Validation métier avec vérification Catalogue
    \item Cache HTTP pour optimiser les performances
    \item Tests PHPUnit validant les cas principaux
\end{itemize}

\end{document}
